%
\hsection{For-Loop / For-Schleife}%
%
%
\begin{question}{What is it?}%
\begin{itemize}%
\item[EN] Explain the syntax \emph{and} purpose/use of an \pythonilIdx{for} statement in \python. %
You can use an example for your explanation.%
\item[DE] Erklären Sie die Syntax \emph{und} den Verwendungszweck eines \pythonil{for}-Statements in \python. %
Sie können ein Beispiel in Ihrer Erklärung verwenden.%
\end{itemize}%
\end{question}%
%
%
\begin{question}{Program Understanding / Programm Verstehen}%
\gitLoadPython{questions:for_01}{}{questions/for_01.py}{}%
\listingPython{questions:for_01}{The program \programUrl{questions:for_01}}%
\begin{itemize}%
\item[EN] Explain program \programUrl{questions:for_01} in \cref{lst:questions:for_01} \emph{line-by-line}, i.e., every single line. %
What output does it produce?
\item[DE] Erklären Sie das Programm \programUrl{questions:for_01} in \cref{lst:questions:for_01} Zeile für Zeile, also jede einzelne Zeile. %
Welche Ausgabe produziert es?%
\end{itemize}%
\end{question}%
%
%
\begin{question}{Program Understanding / Programm Verstehen}%
\gitLoadPython{questions:for_02}{}{questions/for_02.py}{}%
\listingPython{questions:for_02}{The program \programUrl{questions:if_01}}%
\begin{itemize}%
\item[EN] Explain program \programUrl{questions:for_02} in \cref{lst:questions:for_02} \emph{line-by-line}, i.e., every single line. %
What output does it produce?
\item[DE] Erklären Sie das Programm \programUrl{questions:for_02} in \cref{lst:questions:for_02} Zeile für Zeile, also jede einzelne Zeile. %
Welche Ausgabe produziert es?%
\end{itemize}%
\end{question}%
%
%
\endhsection%
%
